% =====================================================================
%  3 Conclusão
% =====================================================================
\section{Conclusão}\label{sec:conclusao}

As atividades práticas da disciplina de Fundamentos de DevOps permitiram construir, de forma
incremental, um laboratório que reproduz em pequena escala o ambiente de uma empresa. O Vagrant,
junto com o VMware, tornou possível criar e recriar sete máquinas virtuais a partir de arquivos de
texto; o Git e o GitHub garantiram o versionamento de todo o trabalho; o Ansible automatizou a
configuração dos servidores de banco de dados e de aplicação e o \textit{deploy} da aplicação Easy
Notes; o Docker empacotou essa aplicação em uma imagem portátil; o SonarQube avaliou a qualidade e a
segurança do código; e o Jenkins encadeou essas etapas em um \textit{pipeline} de CI/CD.

Mais do que aprender comandos, as práticas mostraram os princípios por trás do DevOps:
\textbf{automação} no lugar de tarefas manuais repetitivas, \textbf{Infraestrutura como Código} e
\textbf{versionamento} para ter ambientes reprodutíveis e rastreáveis, \textbf{idempotência} para
executar a automação com segurança quantas vezes for preciso, \textbf{qualidade contínua} com a
análise estática e o Quality Gate, e \textbf{\textit{feedback} rápido} por meio do \textit{pipeline}.
Os problemas enfrentados no caminho -- autenticação SSH entre as VMs, indentação do YAML,
permissões de diretórios, versões do Java, portas e \textit{firewall} -- também fizeram parte do
aprendizado, pois exigiram diagnosticar o ambiente camada por camada, como acontece no dia a dia de
operações.

Outro aprendizado importante foi o cuidado com informações sensíveis: senhas, chaves privadas e
\textit{tokens} não devem aparecer no código versionado nem nos prints, e ferramentas como o Ansible
Vault e as \textit{credentials} do Jenkins existem justamente para isso. Como evolução do laboratório,
seria interessante proteger as variáveis com o Ansible Vault, fazer o Jenkins acionar os
\textit{playbooks} do Ansible, publicar as imagens em um \textit{registry}, adicionar monitoramento
(por exemplo, com Prometheus e Grafana) e levar a mesma estrutura para um ambiente de nuvem. Conclui-se
que as ferramentas estudadas, quando integradas, transformam a entrega de software em um processo
padronizado, rápido e confiável, que é exatamente o que a cultura DevOps propõe.
